% !PRINTED 248
\weberpage{248}
For the function $\gamma_3(\omega)$ we have [by § 54, (15)]:
\[
\gamma_3\left(\frac{\gamma+\delta\omega}{\alpha+\beta\omega}\right) = (-1)^{\alpha\gamma+\beta\gamma+\beta\delta} \gamma_3(\omega),
\]
and in the compositions (8), (10), § 69 we have
\[
\lambda \equiv 1, \quad \alpha \equiv \delta \equiv 0 \ (\mathrm{mod}\ 2).
\]
The $\nu$ quantities
\weq{2}{\gamma_3\left(\frac{c+\partial\omega}{a}\right) \gamma_3(\omega)}
are therefore permuted among themselves, and are accordingly likewise the
roots of a transformation equation of the first level.

In order to form the former of these equations for the simplest case
$n=2$, note the relation [§ 54, (5)]:
\weq{3}{\gamma_2(\omega) = \frac{f(\omega)^{24}-16}{f(\omega)^8} = \frac{f_1(\omega)^{24}+16}{f_1(\omega)^8} = \frac{f_2(\omega)^{24}+16}{f_2(\omega)^8},}
from which, in virtue of
\[
f_1(2\omega) f_2(\omega) = \sqrt{2} \qquad [\S\ 34, (16)]
\]
there follows:
\weq{4}{
\begin{gathered}
\gamma_2(2\omega) = \frac{2^8+f_2(\omega)^{24}}{f_2(\omega)^{16}} \\[4pt]
\gamma_2\left(\frac{\omega}{2}\right) = \frac{2^8+f_1(\omega)^{24}}{f_1(\omega)^{16}} \\[4pt]
\gamma_2\left(\frac{\omega+3}{2}\right) = \frac{2^8-f(\omega)^{24}}{f(\omega)^{16}}. \qquad [\S\ 34, (13)]
\end{gathered}
}

If we denote these three quantities by $x$, $x_0$, $x_1$, then from the
relations (6), (8), § 54 there results:
\[
\begin{aligned}
x + x_0 + x_1 &= \gamma_2(\omega)^2 \\
xx_0 + xx_1 + x_0x_1 &= 495\,\gamma_2(\omega) \\
xx_0x_1 &= -j(\omega) + 2^4\cdot3^3\cdot5^3,
\end{aligned}
\]
so that $x$, $x_0$, $x_1$ are the roots of the equation
\weq{5}{x^3 - \gamma_2(\omega)^2 x^2 + 5\cdot9\cdot11\cdot\gamma_2(\omega) x + j(\omega) - 2^4\cdot3^3\cdot5^3 = 0.}
The equation whose roots are the cubes of the roots of (5) is the
invariant equation for $n=2$, and it can be computed from this without
difficulty.

\begin{center}
\textbf{§ 72. Multiplier Equations of the First Level.}
\end{center}

We shall consider here those multiplier equations whose roots are the
various powers of $P_{11}$ multiplied by powers of
$\gamma_2(\omega)$ and $\gamma_3(\omega)$, and whose coefficients
%CONT>>
